\documentclass[10pt]{article}

\usepackage[utf8]{inputenc}

\usepackage[T1]{fontenc}

\usepackage{amsmath}

\usepackage{amsfonts}

\usepackage{amssymb}

\usepackage[version=4]{mhchem}

\usepackage{stmaryrd}

\usepackage{fixltx2e}

\usepackage{bbold}



\begin{document}

для упрощения принято $x^{\prime}=t^{\prime}=0$. Полученная Г. ф. называется запазды в ю щей, поскольку она обращается в нуль при $t-t^{\prime}<0$. Подставляя Г. Ф. в (3), получим решение неоднородного волнового уравнения в виде



\[

u(x, t)=\left(a^{2} / 4 \pi\right) \int d x^{\prime}\left|x-x^{\prime}\right|^{-1} f\left(x^{\prime}, t-a^{-1}\left|x-x^{\prime}\right|\right),

\]



носящем в электродинамике название з а п а зды в а ю щ е го потенциала.



\begin{enumerate}

  \setcounter{enumi}{5}

  \item Уравнение Клейна - Гордона:

\end{enumerate}



\[

L=\square_{a}+m^{2}

\][



G\_\{0\}(x, t)=(2 \textbackslash pi a)\^{}\{-1\} \textbackslash theta(a t) \textbackslash delta\textbackslash left(a\^{}\{2\} t\textsuperscript{\{2\}-|x|}\{2\}\textbackslash right)-\textbackslash left(m / 4 \textbackslash pi a\^{}\{3\}\textbackslash right) \textbackslash times



][



\textbackslash times \textbackslash theta(a t-|x|)\textbackslash left(t\textsuperscript{\{2\}-a}\{-2\}|x|\textsuperscript{\{2\}\textbackslash right)}\{-1 / 2\} J\_\{1\}\textbackslash left(m \textbackslash sqrt\{t\textsuperscript{\{2\}-a}\{-2\}|x|\^{}\{2\}\}\textbackslash right)

]



при $n=3$, где $J_{1}$ - функция Бесселя. Полученная Г. ф. также называется з а п а зды в а ю щ е й.



Г.ф. играет важную роль также в задачах о спектре дифференциальных операторов. Если самосопряженный оператор $L$ имеет Г. ф., то задача на собственные значения



\[

L u=\lambda u

\]



' эквивалентна интегральному уравнению



\[

u(x)=\lambda \int G\left(x, x^{\prime}\right) u\left(x^{\prime}\right) d x^{\prime},

\]



к к-рому можно применить теорию Фредгольма. Задача $L u=\lambda u$ имеет не более счетного числа собственных значений $\lambda_{1}, \lambda_{2}, \lambda_{3}, \ldots$, все $\lambda_{i}$ вещественны и не имеют конечных точек сгущения. Если комплексное число $\lambda$ не является собственным значением оператора $L$, то можно построить $\Gamma$. Ф. $G\left(x, x^{\prime} ; \lambda\right)$ оператора $L-\lambda I$, где $I-$ единичный оператор. Функция $G\left(x, x^{\prime} ; \lambda\right)$, называемая рез оль вен т ой оператора $L$, является мероморфной функцией параметра $\lambda$, причем ее полюсами служат собственные значения оператора $L$. Таким образом, спектр оператора $L$ можно найти, изучая его резольвенту $G\left(x, x^{\prime} ; \lambda\right)$.



При изучении систем уравнений $L \boldsymbol{u}=\boldsymbol{f}$ роль Г. ф. играют так наз. матр ицы Грина. Они позволяют выразить решение неоднородной краевой задачи для системы в виде интегралов от произведений матрицы Грина на векторы правой части системы. Для подобных задач полезен интеграл Дюамеля. Напр., частное решение неоднородной системы



\[

\boldsymbol{u}^{\prime}=A(x) \boldsymbol{u}+\boldsymbol{F}(x),

\]



где $\boldsymbol{u}$ и $\boldsymbol{F}-n$-компонентные векторы, $A(x)-$ квадратная матрица порядка $n$, записывают в виде



\[

\boldsymbol{u}(x)=\int_{x_{0}}^{x} \boldsymbol{w}(x, s) d s

\]



где $\boldsymbol{w}(x, s)$ - решение однородной системы $\boldsymbol{w}^{\prime}=A(x) \boldsymbol{w}$, $\left.\boldsymbol{w}(x, s)\right|_{x=s}=\boldsymbol{F}(s)$. Матрица $\boldsymbol{A}(x)$ может содержать дифференциальные операторы, поэтому метод применим к уравнениям с частными производными.



Лит.: [1] М орс Ф. М., Фе шбах Г., Методы теоретической физики, пер. с англ., т. 1, М., 1958; [2] Курант Р., Уравнения с частными производными, пер. с англ., М., 1964; [3] Владими ро в В. С., Уравнения математической физики, 4 изд., М., 1981. Л. П. Купцов.



ГРИНОВА ЕМКОСТЬ - см. Емкость множества.



ГРИОЛИ ПРЕЦЕССИЯ - см. Регулярная прецессия.



ГРИФФИТСА ЛЕММА: если последовательность выпуклых дифференцируемых функций сходится поточечно к предельной функции (необходимо выпуклой), то и последовательность производных сходится к производной от предельной функции в точках ее непрерывной дифференцируемости. Более обшая лемма сформулирована М. Фишером (см. [2]): пусть дана последовательность выпуклых вниз функций $\left\{f_{k}(x)\right\}, \quad x \in \mathbf{I} \subset \mathbb{R}^{1}$, к-рая сходится при $k \rightarrow \infty$ к предельной функции $f(x)$; тогда для производных слева, $d^{-} f / d x$, и производных справа, $d^{+} f / d x$, в любой точке $x \in \mathbf{I}$ справедливы неравенства



\[

\frac{d^{-} f(x)}{d x} \leqslant \lim _{k \rightarrow \infty} \inf \frac{d^{-} f_{k}(x)}{d x} \leqslant \lim _{k \rightarrow \infty} \sup \frac{d^{+} f_{k}(x)}{d x} \leqslant \frac{d^{+} f(x)}{d x} .

\]



Лемма Фишера используется в методе аппроксимирующего гамильтониана для вывода неравенств, обобщающих уравнение самосогласования в термодинамич. пределе.



Jum.: [1] Griffith s R. B., «J. Math. Phys.», 1964, v. 5, p. 1215; [2] Fischer M. E., там же, 1965, v. 6, p. 1643-53. Й. Г. Бранков. ГРИФФИТСА НЕРАВЕНСТВА - неравенства для корреляторов классической спиновой системы Изинга. Пусть $\Lambda=\left\{x_{1}, \ldots, x_{V}\right\}-$ конечное множество гиперкубич. решетки $\mathbb{Z}^{v}$ и пусть $\sigma_{x}$ равно 1 или -1 для каждого узла $x \in \Lambda$. Для всякого подмножества $X \subset \Lambda$ положим



\[

\sigma^{X}=\prod_{x \in X} \sigma_{x}

\]



и пусть $J_{X} \geqslant 0$ или $J_{X}=+\infty$, а $J_{\varnothing}=0$. Введем обозначения



\[

\begin{gathered}

U^{\prime}\left(\sigma_{x_{1}}, \ldots, \sigma_{x_{V}}\right)=-\sum_{X \subset \Lambda} J_{X} \sigma^{X}, \\

Z=\sum_{\sigma} \exp \left[-U^{\prime}\left(\sigma_{x_{1}}, \ldots, \sigma_{x_{V}}\right)\right]

\end{gathered}

\]



коррелятор



\[

\left\langle\sigma^{X}\right\rangle=Z^{-1} \sum_{\sigma} \sigma^{X} \exp \left[-U^{\prime}\left(\sigma_{x_{1}}, \ldots, \sigma_{x_{V}}\right)\right] .

\]



Сумма по $\sigma$ производится по всем $2^{V}$ наборам $\sigma_{x_{1}}, \ldots, \sigma_{x_{V}}$.

В этих обозначениях Г. н. имеют вид:



\[

\left\langle\sigma^{X}\right\rangle \geqslant 0 ;\left\langle\sigma^{X} \sigma^{Y}\right\rangle \geqslant\left\langle\sigma^{X} X \sigma^{Y}\right\rangle .

\]



Г. н. используются для доказательства сушествования фазового перехода 1-го рода для решеток Изинга. Г. Н. перенесены на произвольные значения спина, а также обобщены на некоммутативные системы.



Jum.: [1] Griffiths R. B., «J. Math. Phys.», 1967, t. 8, p. 478;



[2] Kelly D. G., Scherman S., там же, 1968, t. 9, p. 466; [3] Р юэль Д., Статистическая механика, пер. с англ., М., 1971.



Д. П. Санкович.



ГРОНУОЛЛА ТЕОРЕМА ПЛОЩАДЕЙ - сМ. Площадей метод.



ГРУБАЯ СИСТЕМА - гладкая динамическая система, обладающая свойством: для любого $\varepsilon>0$ найдется такое $\delta>0$, что при любом ее возмущении, отстоящем от нее в $C^{1}$-метрике не более чем на $\delta$, существует гомеоморфизм фазового пространства, который сдвигает точки не более чем на $\varepsilon$ и переводит траектории невозмушенной системы в траектории возмущенной. Формально определение предполагает заданной нек-рую риманову метрику на фазовом многообразии. Фактически о Г. с. обычно говорят либо когда фазовое многообразие замкнуто, либо когда траектории входят в нек-рую компактную область $G$ с гладкой границей, не касаясь последней, причем возмущение и гомеоморфизм рассматривают только на $G$. Ввиду компактности выбор метрики не играет роли.



Таким образом, при малом (в смысле $C^{1}$ ) возмущении Г. с. получается система, эквивалентная исходной по всем своим топологич. свойствам (однако приведенное определение содержит еще дополнительное требование, чтобы эта эквивалентность осуществлялась посредством гомеоморфизма, близкого к тождественному). Иногда термин «грубость» употребляют в более широком смысле, напр. имея в виду только сохранение при малых возмущениях того или иного свойства системы (в этом случае лучше говорить о грубости данного свойства).



Г. с. были введены А. А. Андроновым и Л. С. Понтрягиным [1]. При малой размерности фазового многообразия (единица для дискретного времени и единица или два для непрерывного) Г.с. допускают простую характеристику в терминах качественных свойств поведения траекторий и образуют открытое всюду плотное множество в пространстве всех динамич. систем, снабженном $C^{1}$-топологией (см. [1], [2]); таким образом, системы с более сложным и более чувствительным к малым возмушениям поведением траекторий можно в этом случае рассматривать как исклю-



11 Математическая физика





\end{document}